\documentclass[border=1cm]{standalone}

\usepackage[T2A]{fontenc}
\usepackage[utf8]{inputenc}
\usepackage[russian]{babel}

\usepackage{tikz}
\usetikzlibrary{math}

\begin{document}

\begin{tikzpicture}

% #1 = x, #2 = y, #3 = размер квадрата, #4 = глубина рекурсии
\newcommand{\SierpinskiCarpet}[4]{%
    \pgfmathtruncatemacro{\Level}{#4}
    \ifnum\Level=0
        \fill (#1,#2) rectangle ++(#3,#3);
    \else
        \pgfmathsetmacro{\Step}{#3/3}
        \pgfmathtruncatemacro{\NextLevel}{\Level-1}
        \foreach \i in {0,1,2}{
            \foreach \j in {0,1,2}{
                % пропускаем центральный квадрат (1,1)
                \ifnum\i=1
                    \ifnum\j=1
                    \else
                        \pgfmathsetmacro{\Xnew}{#1 + \i*\Step}
                        \pgfmathsetmacro{\Ynew}{#2 + \j*\Step}
                        \SierpinskiCarpet{\Xnew}{\Ynew}{\Step}{\NextLevel}
                    \fi
                \else
                    \pgfmathsetmacro{\Xnew}{#1 + \i*\Step}
                    \pgfmathsetmacro{\Ynew}{#2 + \j*\Step}
                    \SierpinskiCarpet{\Xnew}{\Ynew}{\Step}{\NextLevel}
                \fi
            }
        }
    \fi
}

% Несколько итераций ковра Серпинского
\foreach \depth/\x in {0/0,1/5,2/10,3/15}{
    \begin{scope}[shift={(\x,0)}]
        \SierpinskiCarpet{0}{0}{4}{\depth}
        \draw[thick] (0,0) rectangle (4,4);
        \node[below] at (2,-0.4) {Итерация \depth};
    \end{scope}
}

\end{tikzpicture}

\end{document}